% !TeX root = BookN.tex

\title{Complex Potentials in Contact Problems for Prestressed Bodies under Moving Loads}
\titlerunning{Complex Potentials in Contact Problems}

\author{Olexander Guz, Stepan Babych, and Yuriy Glukhov}

\institute{%
	Olexander Guz
	\at \AffIMech
	\\
	\email{guz@inmech.kyiv.ua}
	\and
	 Stepan Babych
	\at \AffIMech
	\\
	\email{babich_sy@ukr.net}
	\and
	Yuriy Glukhov
	\at \AffIMech
	\\
	\email{gluchov.uriy@gmail.com}
}

\maketitle

\graphicspath{{13_GuzBabych/}}


\abstract{This chapter investigates the influence of a protective coating, initial
	stresses, material properties, and surface load velocity on the stress-strain state of a
	compressible elastic half-space. Two models of a layered medium are considered: (1)
	a surface layer represented by concentrated masses, and (2) an elastic plate bonded
	to an elastic half-space. A concentrated force moves along the free surface of the
	protective layer at constant speed. The problems are solved using the method of
	complex potentials, which allows for an exact analytical formulation.
}

%\textbf{Key words:} layered half-space, initial (residual) stresses, moving load,\textbf{ }complex potentials.


\section{Introduction}

In planar dynamic problems for elastic bodies with initial stresses, a special class arises in which the original dynamic formulation can be transformed into a stationary problem in a moving coordinate system that translates at constant velocity. In such cases, the solution can be constructed using complex potentials~\cite{GuzBabych9}. The stress and displacement fields expressed in terms of complex potentials retain the same formal structure in both static and dynamic settings.

In the absence of initial stresses, these complex potentials reduce to well-known representations for dynamic problems in linear elasticity, originally developed by Galin~\cite{GuzBabych8}. For static problems without initial stresses, the type of complex potentials depends on the roots of the characteristic equation: with equal roots, they reduce to the classical Kolosov--Muskhelishvili potentials for isotropic materials~\cite{GuzBabych13}, whereas with unequal roots, they correspond to the Lekhnitskii potentials for orthotropic materials~\cite{GuzBabych12}. When initial stresses are present under static loading, the complex potentials transform into those associated with the linearized plane theory of elasticity with initial stress fields~\cite{GuzBabych9}.

This study addresses plane problems involving the action of a moving surface load on a compressible elastic half-space with both a protective coating and initial stresses. Two models for the two-layer half-space are analyzed and compared:
\begin{enumerate}
    \item a model in which the upper layer (protective coating) is represented by concentrated masses;
    \item a model in which the motion of the surface layer is governed by a system of equations from classical plate theory.
\end{enumerate}

The influence of the protective coating, initial stress field, material mechanical properties, and the motion parameters of the surface load on the stress--strain state of the elastic base is studied using the method of complex potentials.

Related problems have been previously analyzed using the integral Fourier transform method~\cite{GuzBabych2,GuzBabych10}. Investigations accounting for initial stresses can be found in works such as~\cite{GuzBabych3,GuzBabych4,GuzBabych5,GuzBabych6,GuzBabych7,GuzBabych11}.


\section{Complex Representations for Compressible and Incompressible Bodies}

Let us consider the complex representations for compressible bodies, following the approach outlined in~\cite{GuzBabych9}.

Along with the Cartesian coordinates in the initially deformed state $(\xi_{1}, \xi_{2})$, we introduce a Cartesian coordinate system $(y_{1}, y_{2})$ that moves uniformly and rectilinearly along the axis $O\xi_{1}$ with speed $v$ to the right. The transformation between the coordinates $\xi_i$ and $y_i$ is given by
\begin{equation} \label{GuzBabych_1_}
    y_{1} = \xi_{1} - \mathrm{v}t, \quad y_{2} = \xi_{2}.
\end{equation}

The equations of motion in terms of displacements under plane strain for a compressible body occupying the region $\xi_{2} + h \le 0$ take the form~\cite{GuzBabych9}:
\begin{equation} \label{GuzBabych_2_}
    L_{11} u_{1} + L_{12} u_{2} = 0, \quad L_{21} u_{1} + L_{22} u_{2} = 0,
\end{equation}
where
\begin{equation} \label{GuzBabych_3_}
    L_{m\alpha} = \tilde{\omega}_{im\alpha\beta} \frac{\partial^{2}}{\partial y_{i} \partial y_{\beta}}
    - \tilde{\rho} \mathrm{v}^{2} \delta_{m\alpha} \frac{\partial^{2}}{\partial y_{1}^{2}} \quad (i,m,\alpha,\beta = 1,2).
\end{equation}

Using the notation in ~\eqref{GuzBabych_3_} and transformation \eqref{GuzBabych_1_}, the steady-state motion equations for compressible bodies become
\begin{multline} \label{GuzBabych_4_}
    \left( \frac{\partial}{\partial y_{1}} - \frac{1}{\eta_{1}} \frac{\partial}{\partial y_{2}} \right)
    \left( \frac{\partial}{\partial y_{1}} - \frac{1}{\bar{\eta}_{1}} \frac{\partial}{\partial y_{2}} \right) \\
    \times
    \left( \frac{\partial}{\partial y_{1}} - \frac{1}{\eta_{2}} \frac{\partial}{\partial y_{2}} \right)
    \left( \frac{\partial}{\partial y_{1}} - \frac{1}{\bar{\eta}_{2}} \frac{\partial}{\partial y_{2}} \right)
    \chi^{(j)} = 0 \quad (j = 1,2),
\end{multline}
where $\eta_{1}$ and $\eta_{2}$ are the roots of the characteristic equation
\begin{equation} \label{GuzBabych_5_}
    \eta^{4} + 2A \eta^{2} + A_{1} = 0.
\end{equation}

The coefficients $A$ and $A_1$ for a compressible body are given by
\begin{equation} \label{GuzBabych_6_}
    \begin{aligned}
        2A \tilde{\omega}_{2222} \tilde{\omega}_{2112} &=
        \tilde{\omega}_{2222} (\tilde{\omega}_{1111} - \tilde{\rho} \mathrm{v}^{2})
        + \tilde{\omega}_{2112} (\tilde{\omega}_{1221} - \tilde{\rho} \mathrm{v}^{2})
        - (\tilde{\omega}_{1122} + \tilde{\omega}_{1212})^{2}, \\
        2A_{1} \tilde{\omega}_{2222} \tilde{\omega}_{2112} &=
        (\tilde{\omega}_{1111} - \tilde{\rho} \mathrm{v}^{2}) (\tilde{\omega}_{1221} - \tilde{\rho} \mathrm{v}^{2}), \\
        \tilde{\rho} \lambda_{1} \lambda_{2} \lambda_{3} &= \rho,
    \end{aligned}
\end{equation}
where $\rho$ is the density of the material in its natural state, and $\delta_{ij}$ is the Kronecker delta. The tensor components $\tilde{\omega}_{im\alpha\beta}$ are determined by the specific problem~\cite{GuzBabych9}.

We now introduce the complex variables
\begin{equation} \label{GuzBabych_7_}
    z_{j} = y_{1} + \eta_{j} (y_{2} + h), \quad \bar{z}_{j} = y_{1} + \bar{\eta}_{j} (y_{2} + h) \quad (j = 1,2).
\end{equation}

Using \eqref{GuzBabych_7_}, \eqref{GuzBabych_4_} transforms to
\begin{equation} \label{GuzBabych_8_}
    \frac{\partial^{4} \chi^{(j)} }{\partial z_{1} \partial z_{2} \partial \bar{z}_{1} \partial \bar{z}_{2}} = 0 \quad (j = 1,2).
\end{equation}

We now consider two cases depending on the roots of the characteristic equation~\eqref{GuzBabych_5_}.

\medskip\emph{Equal roots.} Let the roots coincide:

\begin{equation} \label{GuzBabych_9_}
    \eta_{1} = \eta_{2} = \eta.
\end{equation}
%
Then the general solution of \eqref{GuzBabych_8_} takes the form
\begin{equation} \label{GuzBabych_10_}
    \chi^{(j)} = \Re\left[ F_{1}^{(j)}(z_{1}) + \bar{z}_{1} F_{2}^{(j)}(z_{1}) \right].
\end{equation}
%
Define the new analytic functions:
\begin{equation} \label{GuzBabych_11_}
    F_{j}^{(1)\prime}(z_{1}) = \eta \phi_{j}(z_{1}), \quad
    F_{j}^{(2)\prime}(z_{1}) = \phi_{j}(z_{1}) \quad (j = 1,2).
\end{equation}

Substituting \eqref{GuzBabych_10_} into the expressions for displacement and stress, and using \eqref{GuzBabych_1_} and \eqref{GuzBabych_11_}, we obtain:
\begin{equation} \label{GuzBabych_12_}
    \begin{aligned}
        \tilde{Q}_{kj} &= \Re \left\{ \gamma_{kj}^{(1)} \left[ \phi_{1}''(z_{1}) + \bar{z}_{1} \phi_{2}''(z_{1}) \right]
        + \gamma_{kj}^{(2)} \phi_{2}'(z_{1}) \right\}, \\
        u_{k} &= \Re \left\{ \gamma_{k}^{(1)} \left[ \phi_{1}'(z_{1}) + \bar{z}_{1} \phi_{2}'(z_{1}) \right]
        + \gamma_{k}^{(2)} \phi_{2}(z_{1}) \right\}.
    \end{aligned}
\end{equation}
%
The coefficients $\gamma_{k}^{(n)}$ and $\gamma_{kj}^{(n)}$ depend on the loading conditions and material properties of the layered medium~\cite{GuzBabych10}.

\medskip\emph{Unequal roots.} Now consider the case when
\begin{equation} \label{GuzBabych_13_}
    \eta_{1} \ne \eta_{2}.
\end{equation}

Then the general solution of \eqref{GuzBabych_8_} is given by
\begin{equation} \label{GuzBabych_14_}
    \chi = 2 \Re \left[ F_{1}(z_{1}) + F_{2}(z_{2}) \right].
\end{equation}

Introduce the analytic functions
\begin{equation} \label{GuzBabych_15_}
    F_{j}''(z_{j}) = \Phi_{j}(z_{j}) \quad (j = 1,2).
\end{equation}

Substituting \eqref{GuzBabych_14_} into the expressions for displacements and stress components, and using Eqs.~\eqref{GuzBabych_1_} and \eqref{GuzBabych_15_}, we arrive at:
\begin{equation} \label{GuzBabych_16_}
    \begin{aligned}
        \tilde{Q}_{ij} &= 2 \Re \left[ \gamma_{ij}^{(1)} \Phi_{1}'(z_{1}) + \gamma_{ij}^{(2)} \Phi_{2}'(z_{2}) \right], \\
        u_{k} &= 2 \Re \left[ \gamma_{k}^{(1)} \Phi_{1}(z_{1}) + \gamma_{k}^{(2)} \Phi_{2}(z_{2}) \right] \quad (i,j,k = 1,2).
    \end{aligned}
\end{equation}

\section{Formulation of Problems for a Two-Layer Prestressed Half-Space Under the Action of a Moving Load}

A two-layer prestressed compressible half-space is considered. The half-space material is assumed to be isotropic and nonlinearly elastic in the unstressed state, with an arbitrary form of the elastic potential. In the case of orthotropic material, the elastically equivalent directions are assumed to coincide with the axes of the chosen coordinate system. The initial stress–strain state of the half-space is assumed to be homogeneous.

The layer and the half-space are defined in the Cartesian coordinate system $(\xi_1, \xi_2, \xi_3)$, introduced in the initially deformed state, and related to the Lagrangian coordinates $(x_1, x_2, x_3)$ of the natural state by the relations $\xi_j = \lambda_j x_j$, where $\lambda_j$ is the elongation $(\lambda_j = \text{const},\, j = 1,2,3)$.

The boundary surfaces of the components are flat and mutually parallel. A linear load $P$ is applied to the free boundary of the layer, moving at constant speed $\mathrm{v}$ over a prolonged period, and assumed to be independent of the coordinate $\xi_3$.

The coordinates of the moving system are introduced according to \eqref{GuzBabych_1_}. It is assumed that the deformation pattern is steady in the moving coordinate system. Furthermore, it is assumed that the stresses induced by the load are much smaller than the initial stresses. This assumption allows the use of the linearized theory of elasticity for initially stressed bodies~\cite{GuzBabych9} to describe the additional stress state caused by the moving load.

\subsection{An Elastic Half-Space With Inhomogeneity in the Form of a Thin Surface Layer}
\label{GuzBabychSec3_1}

A surface layer of thickness $h$ is modeled by concentrated masses with density $\rho_1$. The layer occupies the region $-h \le \xi_2 \le 0$, while the half-space occupies $\xi_2 + h \le 0$. Assuming rigid contact at $y_2 = -h$, the boundary conditions are:
\begin{equation} \label{GuzBabych_17_}
    \tilde{Q}_{21} = P_1(y_1) + h \rho_1 \ddot{u}_1, \quad \tilde{Q}_{22} = P_2(y_1) + h \rho_1 \ddot{u}_2.
\end{equation}

\noindent Here, $P_1(y_1)$ and $P_2(y_1)$ are the tangential and normal stresses on the free surface of the layered half-space.

Under these assumptions, the problem reduces to a plane steady-state formulation governed by the equations of motion in \eqref{GuzBabych_2_} with operators defined in \eqref{GuzBabych_3_}, subject to boundary conditions~\eqref{GuzBabych_17_} and the condition of decay at infinity.

We further present solutions only for the case of unequal roots of \eqref{GuzBabych_5_}, whereas the case of equal roots is detailed in~\cite{GuzBabych10}.

Using complex potentials, the boundary conditions at $y_2 = -h$ can be expressed using \eqref{GuzBabych_17_} and the stress representation in \eqref{GuzBabych_16_} as:
\begin{equation} \label{GuzBabych_18_}
    \begin{aligned}
        2\,\Re \big\{ & \gamma_{2j}^{(1)} \Phi_1(y_1) + \gamma_{2j}^{(2)} \Phi_2(y_1) \\
        & - (-1)^j \rho_1 h v^2 [\gamma_j^{(1)} \Phi_1'(y_1) + \gamma_j^{(2)} \Phi_2'(y_1)] \big\}
        = P_j(y_1) \quad (j = 1,2).
    \end{aligned}
\end{equation}

Thus, in the case of unequal roots of \eqref{GuzBabych_5_}, the solution reduces to solving the system of equations~\eqref{GuzBabych_18_} for the functions $\Phi_j$ ($j=1,2$).

\subsection{Elastic Plate on a Prestressed Half-Space}
\label{GuzBabychSec3_2}

Consider now an elastic layered half-space consisting of an elastic plate of thickness $2h$ and a compressible isotropic half-space with arbitrary elastic potential and initial stresses. The plate occupies the region $-h \le \xi_2 \le h$, and the half-space occupies $\xi_2 + h \le 0$.

The dynamic behavior of the plate is modeled using a system of plate theory equations that incorporate both shear and rotational inertia. In the moving coordinate system defined by \eqref{GuzBabych_1_}, the steady-state equations of motion for the plate under tangential and normal loading are~\cite{GuzBabych10,GuzBabych1}:

\begin{equation} \label{GuzBabych_19_}
    \begin{aligned}
        \theta_1 \frac{d^2 u}{d y_1^2} - \tau &= P_1, \\
        \theta_3 \frac{\partial^2 w}{\partial y_1^2} - 2\kappa G_1 h \frac{\partial \phi}{\partial y_1} - q &= P_2, \\
        \theta_2 \frac{\partial^2 \phi}{\partial y_1^2} + 2\kappa G_1 \left( \frac{\partial w}{\partial y_1} - \phi \right) - \tau &= 0,
    \end{aligned}
\end{equation}
%
where
\[
\theta_1 = 2h\left(\frac{2G_1}{1 - \nu_1} - \rho_1 v^2\right), \;\;
\theta_2 = \frac{2h^2}{3}\left(\frac{2G_1}{1 - \nu_1} - \delta_0 \rho_1 v^2\right), \;\;
\theta_3 = 2h\left(\kappa G_1 - \rho_1 v^2\right).
\]

In \eqref{GuzBabych_19_}, $G_1$, $\nu_1$, and $\rho_1$ are the shear modulus, Poisson's ratio, and density of the plate material; $u$ and $w$ are the displacements of the plate mid-plane at $\xi_2 = 0$; $\delta_0$ is a switch that equals 1 or 0 depending on whether rotational inertia is included; $\phi$ is the rotation angle of the cross-section; and $\kappa$ is the Timoshenko shear coefficient. $P_1$ and $P_2$ are the surface tangential and normal stresses.

Rigid contact between the plate and half-space at $y_2 = -h$ yields:
\begin{equation} \label{GuzBabych_20_}
    \tilde{Q}_{21} = \tau, \quad \tilde{Q}_{22} = q, \quad u_2 = w, \quad u_1 = u + h\phi.
\end{equation}

Thus, we obtain a steady-state problem in the $y_1 O y_2$ plane, consisting of the joint solution of the half-space motion equations~\eqref{GuzBabych_2_}, the plate equations~\eqref{GuzBabych_19_}, and the boundary conditions~\eqref{GuzBabych_20_} under decay conditions at infinity.

\medskip\emph{Unequal roots.} From \eqref{GuzBabych_20_} and \eqref{GuzBabych_16_}, the boundary conditions at $y_2 = -h$ are:
\begin{equation} \label{GuzBabych_21_}
    \begin{aligned}
        2\Re\left[\gamma_{21}^{(1)} \Phi_1'(y_1) + \gamma_{21}^{(2)} \Phi_2'(y_1)\right] &= \tau(y_1), \\
        2\Re\left[\gamma_{22}^{(1)} \Phi_1'(y_1) + \gamma_{22}^{(2)} \Phi_2'(y_1)\right] &= q(y_1), \\
        2\Re\left[\gamma_{2}^{(1)} \Phi_1(y_1) + \gamma_{2}^{(2)} \Phi_2(y_1)\right] &= w(y_1), \\
        2\Re\left[\gamma_{1}^{(1)} \Phi_1(y_1) + \gamma_{1}^{(2)} \Phi_2(y_1)\right] &= u(y_1) + h \phi(y_1).
    \end{aligned}
\end{equation}


Eliminating the unknown functions $u(y_1)$, $w(y_1)$, $q(y_1)$, and $\tau(y_1)$ from \eqref{GuzBabych_19_} and \eqref{GuzBabych_21_}, we obtain a system of equations for the analytic functions $\Phi_i$ $(i=1,2)$ and for the rotation function $\phi$:
\begin{align}
		\notag
       	2\Re \Big\{ -\theta_1 \theta_2& \left[ \gamma_1^{(1)} \Phi_1'''(y_1) + \gamma_1^{(2)} \Phi_2'''(y_1) \right]
        \\ \notag
        + &\left[ \theta_4 \gamma_{21}^{(1)} + \theta_1(\theta_3 - 2\kappa h)\gamma_2^{(1)} \right] \Phi_1''(y_1)  
        \\ \notag
        +&\left[ \theta_4 \gamma_{21}^{(2)} + \theta_1(\theta_3 - 2\kappa h)\gamma_2^{(2)} \right] \Phi_2''(y_1)
        \\ \notag
        - &\theta_1 \left[ \gamma_{22}^{(1)} \Phi_1'(y_1) + \gamma_{22}^{(2)} \Phi_2'(y_1) \right] \Big\}  
        = \theta_1 P_2(y_1) - \theta_2 P_1'(y_1), 
        \\ \notag
        2\Re \Big\{ -\theta_1 \theta_3 &\left[ \gamma_2^{(1)} \Phi_1'''(y_1) + \gamma_2^{(2)} \Phi_2'''(y_1) \right]
        \\ \notag
        + \theta_1& \left[ (2\kappa \gamma_1^{(1)} + \gamma_{22}^{(1)}) \Phi_1''(y_1) + (2\kappa \gamma_1^{(2)} + \gamma_{22}^{(2)}) \Phi_2''(y_1) \right]
          \\ \notag
         - 2\kappa &\left[ \gamma_{21}^{(1)} \Phi_1'(y_1) + \gamma_{21}^{(2)} \Phi_2'(y_1) \right] \Big\}
         = 2\kappa P_1(y_1) - \theta_1 P_2'(y_1), 
		\\ \notag
        \phi(y_1) = \frac{1}{2\kappa h \theta_1} &\Big\{ -\theta_2 P_1(y_1)
        \\ \notag
        + 2\,\Re \Big[ \theta_1 \theta_2 &\left( \gamma_1^{(1)} \Phi_1''(y_1) + \gamma_1^{(2)} \Phi_2''(y_1) \right) 
         \\ \notag
        + & \left( 2\kappa h \theta_1 \gamma_2^{(1)} - \theta_4 \gamma_{21}^{(1)} \right) \Phi_1'(y_1)
        \\ \label{GuzBabych_22_}
        + &\left( 2\kappa h \theta_1 \gamma_2^{(2)} - \theta_4 \gamma_{21}^{(2)} \right) \Phi_2'(y_1) \Big] \Big\}.
\end{align}


Thus, the steady-state response of a two-layer half-space with initial stresses to a moving load in the case of unequal roots of \eqref{GuzBabych_5_} reduces to solving for the analytic functions $\Phi_i$ and the rotation angle $\phi$ using \eqref{GuzBabych_22_}.

\section{Method of Solving Problems in Complex Potentials}

To solve the problems stated in Sections~\ref{GuzBabychSec3_1} and~\ref{GuzBabychSec3_2}, we apply the method of Muskhelishvili~\cite{GuzBabych13}, based on the use of Cauchy integrals for a half-plane. According to~\cite{GuzBabych13}, for an arbitrary holomorphic function $f(z)$ in the lower half-plane $y_{2}+h<0$, which is continuous up to and including the boundary, the following relations hold
\begin{equation}
	\label{GuzBabych_23_}
	\frac{1}{2\pi i}\int_{-\infty}^{+\infty}\frac{f(y_{1})\,dy_{1}}{y_{1}-z}
	=
	-f(z)+\frac{1}{2}a,
	\qquad
	\frac{1}{2\pi i}\int_{-\infty}^{+\infty}\frac{\bar{f}(y_{1})\,dy_{1}}{y_{1}-z}
	=
	-\frac{1}{2}\bar{a},
\end{equation}
where
\[
z = y_{1}+i(y_{2}+h).
\]

In \eqref{GuzBabych_23_}, it is assumed that the function $f(y_{1})$ behaves at large $|y_{1}|$ as
\begin{equation}
	\label{GuzBabych_24_}
	f(y_{1})
	=
	a+o(|y_{1}|^{-\varepsilon})
	=
	f(\infty)+o(|y_{1}|^{-\varepsilon}),
	\qquad
	\varepsilon=\mathrm{const}>0 .
\end{equation}

With these assumptions, we proceed to the study of the problem in the plane $y_{1}Oy_{2}$. The behaviour of the complex potentials $\Phi_{j}(z_{j})$ and $\varphi_{j}(z_{1})$ at infinity is assumed to satisfy the same conditions as in the linear theory of elasticity~\cite{GuzBabych12,GuzBabych13}.

\subsection{Elastic Half-Space With Inhomogeneity in the Form of a Thin Surface Layer}

\medskip\emph{Unequal roots.} Let condition~\eqref{GuzBabych_19_} be satisfied. Applying \eqref{GuzBabych_23_} and~\eqref{GuzBabych_24_} to the system of equations \eqref{GuzBabych_18_}, and assuming rigid contact between the protective layer and the half-space, we obtain the following system of two ordinary differential equations with respect to the functions $\Phi_j(z)$ $(j=1,2)$:
\begin{equation}
	\label{GuzBabych_25_}
	L[\Phi_j(z)] = f_j(z), \qquad j=1,2 .
\end{equation}

For the case of rigid contact, the differential operator $L$ and the right-hand sides have the form
\begin{equation}
	\begin{array}{rcl}
	L &=&
	\begin{aligned}[t]
	 \rho_1^2 h^2 v^4
	&\left(\gamma_1^{(1)}\gamma_2^{(2)}-\gamma_2^{(1)}\gamma_1^{(2)}\right)
	\frac{d^2}{dz^2}
	\\
	+\,\rho_1 h v^2
	\Big[
	&\left(\gamma_1^{(2)}\gamma_{22}^{(1)}+\gamma_2^{(2)}\gamma_{21}^{(1)}\right)
	-
	\left(\gamma_1^{(1)}\gamma_{22}^{(2)}+\gamma_2^{(1)}\gamma_{21}^{(2)}\right)
	\Big]\frac{d}{dz}
	\\
	+
	&\left(
	\gamma_{22}^{(1)}\gamma_{21}^{(2)}-
	\gamma_{22}^{(2)}\gamma_{21}^{(1)}
	\right), 
\end{aligned}
\\
	f_j(z)&=&
	\begin{aligned}[t]
	\frac{(-1)^{j-1}}{2\pi i}
	\Bigg\{
	&\int_{-\infty}^{+\infty}
	\frac{\gamma_{22}^{(m)}P_1(y_1)-\gamma_{21}^{(m)}P_2(y_1)}{y_1-z}\,dy_1
	\\
	+\,\rho_1 h v^2
	&\int_{-\infty}^{+\infty}
	\frac{\gamma_2^{(m)}P_1(y_1)+\gamma_1^{(m)}P_2(y_1)}{(y_1-z)^2}\,dy_1
	\Bigg\},
\end{aligned}
\end{array}
\label{GuzBabych_26_}
\end{equation}
where $j,m=1,2$ and $j\neq m$.

Since at $y_2=-h$ we have $y_1=z=z_1=z_2$, expressions~\eqref{GuzBabych_25_} and~\eqref{GuzBabych_26_} can also be considered in the planes $z_1$ and $z_2$.

Thus, in the case of unequal roots, the problem reduces to solving two ordinary inhomogeneous differential equations of the form~\eqref{GuzBabych_25_} for the functions $\Phi_j(z_j)$ $(j=1,2)$ with the right-hand sides defined in the second equation of \eqref{GuzBabych_26_}. The components of the stress–strain state of the half-space are determined using relations~\eqref{GuzBabych_16_} together with the expressions for the functions $\Phi'_j(z_j)$ $(j=1,2)$.

\subsection{Elastic Plate on a Prestressed Half-Space}

\textit{Unequal roots.} Applying formulas~\eqref{GuzBabych_23_} and~\eqref{GuzBabych_24_} to the first two equations of system~\eqref{GuzBabych_22_}, and taking into account that the stresses vanish at infinity, we obtain, for rigid contact between the plate and the half-space, the following system of two ordinary differential equations with respect to the functions $\Phi'_j(z)$ $(j=1,2)$:
\begin{equation}
	\label{GuzBabych_27_}
	L[\Phi'_j(z)] = f_j(z), \qquad j=1,2 .
\end{equation}

The differential operator $L$ and the right-hand side of \eqref{GuzBabych_27_} have the form
\begin{equation}\label{GuzBabych_28_}
	\begin{array}{l}
		\displaystyle
		L =
		\begin{aligned}[t]
			&\theta_1\theta_2\theta_3
			\Big(\gamma_2^{(1)}\gamma_1^{(2)}-\gamma_2^{(2)}\gamma_1^{(1)}\Big)
			\frac{d^4}{dz^4}
			\\
			&+
			\Big[
			\theta_1\theta_2
			\Big(\gamma_1^{(1)}\gamma_{22}^{(2)}-\gamma_1^{(2)}\gamma_{22}^{(1)}\Big)
			+
			\theta_3\theta_4
			\Big(\gamma_{21}^{(1)}\gamma_2^{(2)}-\gamma_2^{(1)}\gamma_{21}^{(2)}\Big)
			\Big]
			\frac{d^3}{dz^3}
			\\
			&+
			\Big\{
			2\kappa\theta_1
			\Big[
			h\Big(\gamma_1^{(1)}\gamma_{21}^{(2)}-\gamma_{21}^{(1)}\gamma_1^{(2)}
			+\gamma_2^{(1)}\gamma_{22}^{(2)}-\gamma_{22}^{(1)}\gamma_2^{(2)}\Big)
			\\
			&\qquad
			+
			(\theta_3-2\kappa h)
			\Big(\gamma_1^{(1)}\gamma_2^{(2)}-\gamma_1^{(2)}\gamma_2^{(1)}\Big)
			\Big]
			+
			\theta_4
			\Big(\gamma_{22}^{(1)}\gamma_{21}^{(2)}-\gamma_{22}^{(2)}\gamma_{21}^{(1)}\Big)
			\Big\}
			\frac{d^2}{dz^2}
			\\
			&+
			2\kappa
			\Big[
			(\theta_3-2\kappa h)
			\Big(\gamma_2^{(1)}\gamma_{21}^{(2)}-\gamma_2^{(2)}\gamma_{21}^{(1)}\Big)
			+
			\theta_1
			\Big(\gamma_{22}^{(1)}\gamma_1^{(2)}-\gamma_1^{(1)}\gamma_{22}^{(2)}\Big)
			\Big]
			\frac{d}{dz}
			\\
			&+
			2\kappa
			\Big(\gamma_{22}^{(2)}\gamma_{21}^{(1)}-\gamma_{22}^{(1)}\gamma_{21}^{(2)}\Big),
		\end{aligned}
		\\[1.2em]
		\displaystyle
		f_j(z) =
		\begin{aligned}[t]
			&\frac{(-1)^j}{2\theta_1\pi i}
			\int_{-\infty}^{+\infty}
			\frac{1}{(y_1-z)^3}
			\Bigg\{
			\Big[
			2\kappa\gamma_{21}^{(k)}(y_1-z)^2
			\\
			&\qquad
			-\theta_1(2\kappa\gamma_1^{(k)}+\gamma_{22}^{(k)})(y_1-z)
			+2\theta_1\theta_3\gamma_2^{(k)}
			\Big]
			[\theta_1 P_2(y_1)-\theta_2 P_1'(y_1)]
			\\
			&\qquad
			+
			\Big[
			\theta_1\gamma_{22}^{(k)}(y_1-z)^2
			+
			(\theta_4\gamma_{21}^{(k)}
			+\theta_1(\theta_3-2\kappa h)\gamma_2^{(k)})(y_1-z)
			\\
			&\qquad\qquad
			-2\theta_1\theta_2\gamma_1^{(k)}
			\Big]
			[2\kappa P_1(y_1)-\theta_1 P_2'(y_1)]
			\Bigg\}
			\,dy_1 ,
		\end{aligned}
	\end{array}
\end{equation}
where $j,k=1,2$ and $j\ne k$.

Since at $y_2=-h$ we have $y_1=z=z_1=z_2$, expressions~\eqref{GuzBabych_27_} and~\eqref{GuzBabych_28_} can also be considered in the planes $z_1$ and $z_2$.

Thus, in the case of unequal roots, the problem of the steady-state response of a two-layer elastic half-space with initial stresses reduces to solving two ordinary inhomogeneous differential equations~\eqref{GuzBabych_27_} for the functions $\Phi'_j(z_j)$. The components of the stress--strain state are determined using relations~\eqref{GuzBabych_16_}, the expressions for $\Phi'_j(z_j)$, and the function $\phi(y_1)$ (see the third equation in system~\eqref{GuzBabych_22_}).

The critical velocities of the moving load are determined from the condition that the characteristic equation of differential equations~\eqref{GuzBabych_27_} possesses real positive multiple roots.

A comparative analysis of the formulas obtained in~\cite{GuzBabych2,GuzBabych10} shows that the equation $\Delta(k)=0$ has the same roots as the characteristic equation corresponding to differential equations~\eqref{GuzBabych_27_}.

Thus, the application of the complex potential method yields results consistent with those obtained by the method of integral Fourier transforms~\cite{GuzBabych2,GuzBabych10}.

\section{Numerical Studies}

Some results illustrating the influence of initial stresses and the mechanical characteristics of the surface layer and the half-space on the components of the stress--strain state of the half-space are presented below. In the numerical analysis, a compressible half-space with a harmonic-type elastic potential was considered~\cite{GuzBabych9,GuzBabych10}.

It is assumed that the initial deformed state is plane, i.e., $\lambda_3 = 1$, and that no surface load is applied in the initial state, $S_{0}^{22} = 0$.

The results of the calculations are illustrated for the case of a concentrated linear load whose normal and tangential components are given by
\[
P_{1} = P\,\delta(y_{1})\cos\alpha,
\qquad
P_{2} = P\,\delta(y_{1})\sin\alpha,
\]
where $\alpha$ is the angle between the direction of the load and the axis $Oy_{1}$.

Let us analyze how the initial stresses in the base influence the characteristics of the stress--strain state for different load velocities.

Figures~\ref{GuzBabychFig1}--\ref{GuzBabychFig3} show the distribution of the generalised stress $\tilde{Q}_{22}$ at the depth
\[
y_{2} = -\frac{2h}{\lambda_{2}}
\]
for rigid contact between the protective layer and the half-space (a: the top layer, i.e., the protective coating, is modelled by concentrated masses; b: an elastic plate on an elastic half-space) and for different velocities of the surface load.

\begin{figure}[b]
	\centering
	\begin{subfigure}{0.48\textwidth}
		\centering
		\includegraphics[width=\linewidth]{13_GuzBabych/GuzBabychFig1a.png}
		\caption{Top layer (protective coating) modelled by concentrated masses.}
	\end{subfigure}\hfill
	\begin{subfigure}{0.48\textwidth}
		\centering
		\includegraphics[width=\linewidth]{13_GuzBabych/GuzBabychFig1b.png}
		\caption{Elastic plate on an elastic half-space.\newline}
	\end{subfigure}
	\caption{Distribution of the generalised stress $\tilde{Q}_{22}$ at the depth $y_2=-2h/\lambda_2$ for subsonic load velocity ($v<c_{12}$).}
	\label{GuzBabychFig1}
\end{figure}

\begin{figure}[t]
	\centering
	\begin{subfigure}{0.48\textwidth}
		\centering
		\includegraphics[width=\linewidth]{13_GuzBabych/GuzBabychFig2a.png}
		\caption{Top layer (protective coating) modelled by concentrated masses.}
	\end{subfigure}\hfill
	\begin{subfigure}{0.48\textwidth}
		\centering
		\includegraphics[width=\linewidth]{13_GuzBabych/GuzBabychFig2b.png}
		\caption{Elastic plate on an elastic half-space.\newline}
	\end{subfigure}
	\caption{Distribution of the generalised stress $\tilde{Q}_{22}$ at the depth $y_2=-2h/\lambda_2$ for transonic load velocity ($c_{12}<v<c_{11}$).}
	\label{GuzBabychFig2}
\end{figure}

\begin{figure}[t]
	\centering
	\begin{subfigure}{0.48\textwidth}
		\centering
		\includegraphics[width=\linewidth]{13_GuzBabych/GuzBabychFig3a.png}
		\caption{Top layer (protective coating) modelled by concentrated masses.}
	\end{subfigure}\hfill
	\begin{subfigure}{0.48\textwidth}
		\centering
		\includegraphics[width=\linewidth]{13_GuzBabych/GuzBabychFig3b.png}
		\caption{Elastic plate on an elastic half-space.\newline}
	\end{subfigure}
	\caption{Distribution of the generalised stress $\tilde{Q}_{22}$ at the depth $y_2=-2h/\lambda_2$ for supersonic load velocity ($v>c_{11}$).}
	\label{GuzBabychFig3}
\end{figure}


The calculations correspond to three characteristic regimes of load motion: subsonic ($v<c_{12}$, Fig.~\ref{GuzBabychFig1}), transonic ($c_{12}<v<c_{11}$, Fig.~\ref{GuzBabychFig2}), and supersonic ($v>c_{11}$, Fig.~\ref{GuzBabychFig3}). Here $c_{11}$ and $c_{12}$ denote the propagation velocities in the $Oy_{1}$ direction of the longitudinally and transversely polarised waves, respectively, in an unconfined body with initial stresses,
\[
\tilde{\rho} c_{ij}^{2} = \tilde{\omega}_{ijji}
\quad (i,j = 1,2) ,
\]
see~\cite{GuzBabych9,GuzBabych10}.


Since the perturbations produced by the moving load are assumed to be small, only subcritical load velocities $v<c_{12}$ were considered in the analysis~\cite{GuzBabych9,GuzBabych10}.

Curves 1--5 in Figs.~\ref{GuzBabychFig1} and~\ref{GuzBabychFig3} correspond to the following values of the initial elongation:
\[
\lambda_{1} = 0.8,\; 0.9,\; 1.0,\; 1.1,\; 1.2 .
\]
The graphs represent the behaviour of points of the half-space that coincide in the natural state.

For the diagrams shown in Figs.~\ref{GuzBabychFig1}--\ref{GuzBabychFig3}, the calculations were carried out for the following parameter values:
\begin{align*}
&P=\mu, \qquad
\frac{\rho}{\rho_{1}}=0.25, \qquad
\frac{\mu}{G_{1}}=0.25,
\\
&\nu=0.3, \qquad
\nu_{1}=0.25, \qquad
\delta_{0}=1, \qquad
\kappa=0.845, \qquad
\alpha=\frac{\pi}{2}.
\end{align*}
Here $\nu$ denotes the Poisson ratio of the half-space material.

Figure~\ref{GuzBabychFig1} shows the distribution of the generalised stress $\tilde{Q}_{22}$ in the half-space at the depth $y_{2}=-2h/\lambda_{2}$ for the subsonic velocity $v^{2}=0.1\,c_{s}^{2}$; Fig.~\ref{GuzBabychFig2} corresponds to the transonic velocity $v^{2}=2\,c_{s}^{2}$; and Fig.~\ref{GuzBabychFig3} corresponds to the supersonic velocity $v^{2}=6\,c_{s}^{2}$.


The analysis of the numerical results leads to the following conclusions. The dependencies of the stress–strain state parameters are qualitatively similar for the considered layered-medium models.


The values of the parameters characterising the stress--strain state of the base depend on the coordinates of the observation point, the initial stresses, the velocity of the moving load, and the mechanical properties of the layered medium.

For subsonic loading speeds, the graphs of the quantities describing the stress--strain state are symmetric with respect to the point of load application (Fig.~\ref{GuzBabychFig1}). The initial stresses have the greatest influence on the stress--strain state of the base in the vicinity of the load application point. The attenuation of stresses with distance from this point is slower under compression than under tension. At the same time, there exist regions of the half-space where the stresses and displacement velocities depend only weakly on the initial deformations.

As the load velocity increases, both the stresses and the displacement velocities in the half-space decrease. At higher velocities, the influence of the initial stresses also becomes weaker. For a stiffer surface layer, the influence of both the load velocity and the initial stresses is reduced.

For the transonic loading regime (Fig.~\ref{GuzBabychFig2}), the distributions of the considered quantities become asymmetric with respect to the load application point. In this case, the forward-propagating wave attenuates much faster than the backward-propagating wave.

The numerical results show that the presence of initial stresses significantly affects the distributions of stresses and displacement velocities in the half-space. This influence depends on the position of the observation point relative to the load application point. The effect of initial stresses becomes more pronounced as the load velocity increases, especially in the case of pre-compression. As in the case of subsonic velocities ($v<c_{12}$), there exist regions of the half-space where the stresses and displacement velocities depend only weakly on the initial deformations.

From Fig.~\ref{GuzBabychFig3} (supersonic regime, $v>c_{11}$), it can be seen that the asymmetry of the distributions increases with velocity, and the forward wave attenuates much faster. However, it does not disappear completely, which is likely due to the layered structure of the medium. A qualitatively similar behaviour was observed in~\cite{GuzBabych2,GuzBabych10}.

The dependence of the stress--strain state parameters on the initial stresses has the same qualitative character as in the case $c_{12}<v<c_{11}$.

For supersonic loading speeds, the diagrams of the considered quantities remain asymmetric with respect to the load application point. In this regime, the forward wave attenuates much faster than the backward wave, but does not vanish completely because of the presence of the surface layer.

The stresses and displacement velocities in the half-space depend significantly on the initial stresses, and the specific form of this dependence is determined by the position of the observation point relative to the load application point.

As in the case of subsonic velocities, there exist regions of the half-space where the stresses and displacement velocities are almost independent of the initial deformations.

At subsonic velocities, the amplitudes of the considered parameters are larger for the first model of the layered medium (concentrated masses on the elastic half-space). At supersonic velocities, this difference becomes insignificant.

\section{Conclusions}

In this paper, the influence of the protective coating, initial stresses, mechanical characteristics of materials, and parameters of surface load movement on the stress-strain state of an elastic basic is investigated using the method of complex potentials.

Two models of a layered compressible half-space are considered and compared: 1) an elastic plate on an elastic half-space; 2) the top layer (protective coating) is modelled by concentrated masses. The concentrated force moves along the free surface of the protective layer with a constant speed.

To solve the problems, we use the Muskhelishvili method based on Cauchy-type integrals for a half-plane. In this case, the problem is reduced to solving systems of ordinary non-homogeneous linear differential equations with constant coefficients with respect to unknown analytical functions.

The analytical results are presented in a general form for compressible materials with an arbitrary elastic potential, for cases of unequal and equal roots of the characteristic equations and for any load speeds.

For the numerical analysis, a compressible material with the harmonic type potential were considered. The calculations were performed within the framework of the theory of finite initial deformations. The influence of the moving load, initial stresses of mechanical parameters of the layered base elements on the main characteristics of its stress-strain state is investigated.

Applying the method of complex potentials, we obtain results analogous to those obtained by the method of integral Fourier transforms.


\bibliographystyle{unsrtnat}
\bibliography{Chapter13}





